\documentclass[10pt,a4paper]{amsart}
\usepackage{amsmath,amssymb,amsthm,graphicx}
% Retained minimal preamble for the Stacks Project chapter more-morphisms.tex (More on Morphisms).
% Verbatim excerpt of the chapter preamble: Xy-pic setup (native preamble lines 14-18) and the native
% theorem declarations with the mathematical macros (native preamble lines 154-217). Copied mechanically;
% nothing added, removed or redefined. The wrapper class and the amsmath/amssymb/amsthm/graphicx base
% packages are supplied by the assembler, so the class/externaldocument part of the native preamble is omitted.
% For commutative diagrams we use Xy-pic
\usepackage[all]{xy}

% We use 2cell for 2-commutative diagrams.
\xyoption{2cell}
% Theorem environments.
%
\theoremstyle{plain}
\newtheorem{theorem}[subsection]{Theorem}
\newtheorem{proposition}[subsection]{Proposition}
\newtheorem{lemma}[subsection]{Lemma}

\theoremstyle{definition}
\newtheorem{definition}[subsection]{Definition}
\newtheorem{example}[subsection]{Example}
\newtheorem{exercise}[subsection]{Exercise}
\newtheorem{situation}[subsection]{Situation}

\theoremstyle{remark}
\newtheorem{remark}[subsection]{Remark}
\newtheorem{remarks}[subsection]{Remarks}

\numberwithin{equation}{subsection}

% Macros
%
\def\lim{\mathop{\mathrm{lim}}\nolimits}
\def\colim{\mathop{\mathrm{colim}}\nolimits}
\def\Spec{\mathop{\mathrm{Spec}}}
\def\Hom{\mathop{\mathrm{Hom}}\nolimits}
\def\Ext{\mathop{\mathrm{Ext}}\nolimits}
\def\SheafHom{\mathop{\mathcal{H}\!\mathit{om}}\nolimits}
\def\SheafExt{\mathop{\mathcal{E}\!\mathit{xt}}\nolimits}
\def\Sch{\mathit{Sch}}
\def\Mor{\mathop{\mathrm{Mor}}\nolimits}
\def\Ob{\mathop{\mathrm{Ob}}\nolimits}
\def\Sh{\mathop{\mathit{Sh}}\nolimits}
\def\NL{\mathop{N\!L}\nolimits}
\def\CH{\mathop{\mathrm{CH}}\nolimits}
\def\proetale{{pro\text{-}\acute{e}tale}}
\def\etale{{\acute{e}tale}}
\def\QCoh{\mathit{QCoh}}
\def\Ker{\mathop{\mathrm{Ker}}}
\def\Im{\mathop{\mathrm{Im}}}
\def\Coker{\mathop{\mathrm{Coker}}}
\def\Coim{\mathop{\mathrm{Coim}}}

% Boxtimes
%
\DeclareMathSymbol{\boxtimes}{\mathbin}{AMSa}{"02}

%
% Macros for moduli stacks/spaces
%
\def\QCohstack{\mathcal{QC}\!\mathit{oh}}
\def\Cohstack{\mathcal{C}\!\mathit{oh}}
\def\Spacesstack{\mathcal{S}\!\mathit{paces}}
\def\Quotfunctor{\mathrm{Quot}}
\def\Hilbfunctor{\mathrm{Hilb}}
\def\Curvesstack{\mathcal{C}\!\mathit{urves}}
\def\Polarizedstack{\mathcal{P}\!\mathit{olarized}}
\def\Complexesstack{\mathcal{C}\!\mathit{omplexes}}
% \Pic is the operator that assigns to X its picard group, usage \Pic(X)
% \Picardstack_{X/B} denotes the Picard stack of X over B
% \Picardfunctor_{X/B} denotes the Picard functor of X over B
\def\Pic{\mathop{\mathrm{Pic}}\nolimits}
\def\Picardstack{\mathcal{P}\!\mathit{ic}}
\def\Picardfunctor{\mathrm{Pic}}
\def\Deformationcategory{\mathcal{D}\!\mathit{ef}}
\newenvironment{reference}{}{}
\newenvironment{slogan}{}{}
\newenvironment{history}{}{}
\usepackage{cleveref}
\crefname{originallemma}{original lemma}{original lemmas}
\Crefname{originallemma}{Original lemma}{Original lemmas}
\makeatletter
\newlabel{algebra-lemma-relative-global-complete-intersection-conormal}{{00SV}{source}{}{sourcepointer-1}{}}
\newlabel{algebra-lemma-relative-global-complete-intersection-conormal@cref}{{[originallemma][0][]00SV}{[1][0][]source}{}{}{}}
\makeatother
\crefname{originallemma}{original lemma}{original lemmas}
\Crefname{originallemma}{Original lemma}{Original lemmas}
\makeatletter
\newlabel{divisors-definition-relative-H1-regular-immersion}{{063S}{source}{}{sourcepointer-2}{}}
\newlabel{divisors-definition-relative-H1-regular-immersion@cref}{{[originallemma][0][]063S}{[1][0][]source}{}{}{}}
\makeatother
\crefname{originallemma}{original lemma}{original lemmas}
\Crefname{originallemma}{Original lemma}{Original lemmas}
\makeatletter
\newlabel{divisors-lemma-relative-regular-immersion-flat-in-neighbourhood}{{063U}{source}{}{sourcepointer-3}{}}
\newlabel{divisors-lemma-relative-regular-immersion-flat-in-neighbourhood@cref}{{[originallemma][0][]063U}{[1][0][]source}{}{}{}}
\makeatother
\crefname{originallemma}{original lemma}{original lemmas}
\Crefname{originallemma}{Original lemma}{Original lemmas}
\makeatletter
\newlabel{lemma-lci-properties}{{069H}{source}{}{sourcepointer-4}{}}
\newlabel{lemma-lci-properties@cref}{{[originallemma][0][]069H}{[1][0][]source}{}{}{}}
\makeatother
\crefname{originallemma}{original lemma}{original lemmas}
\Crefname{originallemma}{Original lemma}{Original lemmas}
\makeatletter
\newlabel{more-algebra-lemma-homotopy-koszul}{{0663}{source}{}{sourcepointer-5}{}}
\newlabel{more-algebra-lemma-homotopy-koszul@cref}{{[originallemma][0][]0663}{[1][0][]source}{}{}{}}
\makeatother
\crefname{originallemma}{original lemma}{original lemmas}
\Crefname{originallemma}{Original lemma}{Original lemmas}
\makeatletter
\newlabel{more-algebra-lemma-regular-koszul-regular}{{062F}{source}{}{sourcepointer-6}{}}
\newlabel{more-algebra-lemma-regular-koszul-regular@cref}{{[originallemma][0][]062F}{[1][0][]source}{}{}{}}
\makeatother
\crefname{originallemma}{original lemma}{original lemmas}
\Crefname{originallemma}{Original lemma}{Original lemmas}
\makeatletter
\newlabel{more-algebra-lemma-syntomic-lci}{{07D3}{source}{}{sourcepointer-7}{}}
\newlabel{more-algebra-lemma-syntomic-lci@cref}{{[originallemma][0][]07D3}{[1][0][]source}{}{}{}}
\makeatother
\crefname{originallemma}{original lemma}{original lemmas}
\Crefname{originallemma}{Original lemma}{Original lemmas}
\makeatletter
\newlabel{morphisms-lemma-syntomic-locally-standard-syntomic}{{01UE}{source}{}{sourcepointer-8}{}}
\newlabel{morphisms-lemma-syntomic-locally-standard-syntomic@cref}{{[originallemma][0][]01UE}{[1][0][]source}{}{}{}}
\makeatother
\title{A morphism of schemes is flat and local complete intersection exactly when it is syntomic}
\author{The Stacks Project Authors}
\date{stacks-project:more-morphisms.tex:local-snapshot-20261011; GNU FDL 1.2 or later}
\begin{document}
\maketitle
\noindent\textit{Editorial note.} Statement span is lines 18244-18255 of sources/stacks-more-morphisms/pinned-original-source.tex; proof span is lines 18257-18303, the complete author proof (45 body lines).\par\medskip
\noindent\textit{Editorial note.} Navigation stubs: 8 cited-but-not-included labels are supplied as non-load-bearing pointers with their pinned tags where the snapshot tag index has them.\par\medskip
\noindent\textit{Editorial note.} extra\_packages is empty on purpose: the retained preamble already loads xy with options and a second declaration clashes in the assembler wrapper.\par\medskip
\noindent\textit{Editorial note.} Text is retained verbatim, including the author's own typos and placeholders.\par\medskip
\noindent\textit{Editorial note.} Counted A side only (B-side proof incomplete per ruling). First pillar of the L16201 descent proof.\par\medskip
\noindent\textit{Allowed standard prerequisites (editorial).}\begin{itemize}
\item Explicit prerequisite boundary. More on Algebra, lemma syntomic-lci; Algebra, lemma relative-global-complete-intersection-conormal; More on Algebra, lemmas regular-koszul-regular and homotopy-koszul; Divisors, lemma relative-regular-immersion-flat-in-neighbourhood; chapter lemma lci-properties.
\end{itemize}\medskip
\ifdefined\hypertarget\hypertarget{sourcepointer-1}{}\fi
\noindent\textit{Source reference.} Original lemma 00SV: sources/stacks-more-morphisms/pinned-original-source.tex, source lines 18303--18303. Cited by the retained statement or proof but not part of the collected spans; supplied only as a navigation pointer. The argument does not depend on its content.\par\medskip
\ifdefined\hypertarget\hypertarget{sourcepointer-2}{}\fi
\noindent\textit{Source reference.} Original lemma 063S: sources/stacks-more-morphisms/pinned-original-source.tex, source lines 18303--18303. Cited by the retained statement or proof but not part of the collected spans; supplied only as a navigation pointer. The argument does not depend on its content.\par\medskip
\ifdefined\hypertarget\hypertarget{sourcepointer-3}{}\fi
\noindent\textit{Source reference.} Original lemma 063U: sources/stacks-more-morphisms/pinned-original-source.tex, source lines 18303--18303. Cited by the retained statement or proof but not part of the collected spans; supplied only as a navigation pointer. The argument does not depend on its content.\par\medskip
\ifdefined\hypertarget\hypertarget{sourcepointer-4}{}\fi
\noindent\textit{Source reference.} Original lemma 069H: sources/stacks-more-morphisms/pinned-original-source.tex, source lines 18150--18157. Cited by the retained statement or proof but not part of the collected spans; supplied only as a navigation pointer. The argument does not depend on its content.\par\medskip
\ifdefined\hypertarget\hypertarget{sourcepointer-5}{}\fi
\noindent\textit{Source reference.} Original lemma 0663: sources/stacks-more-morphisms/pinned-original-source.tex, source lines 18303--18303. Cited by the retained statement or proof but not part of the collected spans; supplied only as a navigation pointer. The argument does not depend on its content.\par\medskip
\ifdefined\hypertarget\hypertarget{sourcepointer-6}{}\fi
\noindent\textit{Source reference.} Original lemma 062F: sources/stacks-more-morphisms/pinned-original-source.tex, source lines 18303--18303. Cited by the retained statement or proof but not part of the collected spans; supplied only as a navigation pointer. The argument does not depend on its content.\par\medskip
\ifdefined\hypertarget\hypertarget{sourcepointer-7}{}\fi
\noindent\textit{Source reference.} Original lemma 07D3: sources/stacks-more-morphisms/pinned-original-source.tex, source lines 18303--18303. Cited by the retained statement or proof but not part of the collected spans; supplied only as a navigation pointer. The argument does not depend on its content.\par\medskip
\ifdefined\hypertarget\hypertarget{sourcepointer-8}{}\fi
\noindent\textit{Source reference.} Original lemma 01UE: sources/stacks-more-morphisms/pinned-original-source.tex, source lines 18303--18303. Cited by the retained statement or proof but not part of the collected spans; supplied only as a navigation pointer. The argument does not depend on its content.\par\medskip
\begin{lemma}
\label{lemma-flat-lci}
\begin{slogan}
A morphism is flat and lci if and only if it is syntomic.
\end{slogan}
Let $f : X \to S$ be a morphism of schemes.
The following are equivalent
\begin{enumerate}
\item $f$ is flat and a local complete intersection morphism, and
\item $f$ is syntomic.
\end{enumerate}
\end{lemma}

\begin{proof}
Working affine locally this is
More on Algebra, Lemma \ref{more-algebra-lemma-syntomic-lci}.
We also give a more geometric proof.

\medskip\noindent
Assume (2). By
Morphisms, Lemma \ref{morphisms-lemma-syntomic-locally-standard-syntomic}
for every point $x$ of $X$ there exist affine open neighbourhoods
$U$ of $x$ and $V$ of $f(x)$ such that $f|_U : U \to V$ is standard syntomic.
This means that $U = \Spec(R[x_1, \ldots, x_n]/(f_1, \ldots, f_c))
\to V = \Spec(R)$ where $R[x_1, \ldots, x_n]/(f_1, \ldots, f_c)$ is a
relative global complete intersection over $R$. By
Algebra,
Lemma \ref{algebra-lemma-relative-global-complete-intersection-conormal}
the sequence $f_1, \ldots, f_c$ is a regular sequence in each local
ring $R[x_1, \ldots, x_n]_{\mathfrak q}$ for every prime
$\mathfrak q \supset (f_1, \ldots, f_c)$. Consider the Koszul complex
$K_\bullet = K_\bullet(R[x_1, \ldots, x_n], f_1, \ldots, f_c)$
with homology groups $H_i = H_i(K_\bullet)$. By
More on Algebra, Lemma \ref{more-algebra-lemma-regular-koszul-regular}
we see that $(H_i)_{\mathfrak q} = 0$, $i > 0$ for every $\mathfrak q$
as above. On the other hand, by
More on Algebra, Lemma \ref{more-algebra-lemma-homotopy-koszul}
we see that $H_i$ is annihilated by $(f_1, \ldots, f_c)$. Hence we
see that $H_i = 0$, $i > 0$ and $f_1, \ldots, f_c$ is a Koszul-regular
sequence. This proves that $U \to V$ factors as a Koszul-regular
immersion $U \to \mathbf{A}^n_V$ followed by a smooth morphism as desired.

\medskip\noindent
Assume (1). Then $f$ is a flat and locally of finite presentation
(Lemma \ref{lemma-lci-properties}).
Hence, according to
Morphisms, Lemma \ref{morphisms-lemma-syntomic-locally-standard-syntomic}
it suffices to show that the local rings $\mathcal{O}_{X_s, x}$
are local complete intersection rings. Choose, locally on $X$, a factorization
$f = \pi \circ i$ for some Koszul-regular immersion $i : X \to P$
and smooth morphism $\pi : P \to S$. Note that $X \to P$ is
a relative quasi-regular immersion over $S$, see
Divisors, Definition \ref{divisors-definition-relative-H1-regular-immersion}.
Hence according to
Divisors,
Lemma \ref{divisors-lemma-relative-regular-immersion-flat-in-neighbourhood}
we see that $X \to P$ is a regular immersion and the same remains true
after any base change. Thus each fibre is a regular immersion, whence
all the local rings of all the fibres of $X$ are local complete intersections.
\end{proof}

\end{document}
